% =========================================================
% PHẦN 5: BÀI HỌC KINH NGHIỆM, HƯỚNG PHÁT TRIỂN VÀ GENAI (CHƯƠNG 5-6)
% =========================================================
\section{Bài học Kinh nghiệm và Khai báo AI}

\begin{frame}
\frametitle{Bài học Kinh nghiệm (Lessons Learned) và Hướng Phát triển}
% TODO: Bắt buộc theo PDF trang 4 ("contains lessons learned for the team")
\begin{itemize}
    \item \textbf{Bài học kinh nghiệm:} Chuẩn hóa hợp đồng giao tiếp (\texttt{API\_CONTRACTS.md}) và ranh giới sở hữu dữ liệu từ sớm; kiểm soát tranh chấp đồng thời ở tầng giao dịch; duy trì tính nhất quán giữa đặc tả, sơ đồ UML và mã nguồn.
    \item \textbf{Hướng phát triển:} Hoàn thiện kết nối Live UI cho các màn hình quản trị (\texttt{UC08/09/11}, Parking/Charging) và tích hợp giao thức thiết bị IoT thực tế (OCPP/MQTT).
\end{itemize}
\end{frame}

\begin{frame}
\frametitle{Khai báo Minh bạch Generative AI (Tools, Scope, Contribution)}
% TODO: Bắt buộc theo PDF trang 5-6 và PROJECT_SPECIFICATION.md (§8)
\begin{itemize}
    \item \textbf{Công cụ và Phạm vi:} Hỗ trợ kiểm tra quy chuẩn UML, gợi ý khung mã kiểm thử và chuẩn hóa trình bày tài liệu LaTeX/Markdown.
    \item \textbf{Kiểm chứng và Đóng góp nguyên bản:} Nhóm trực tiếp phân tích bài toán, thiết kế kiến trúc, xử lý các lỗi logic đồng thời và thực thi kiểm thử xác nhận tính đúng đắn.
\end{itemize}
\end{frame}

\begin{frame}
\frametitle{Cảm ơn và Hỏi đáp (Q\&A)}
\begin{center}
    {\Large \textbf{Cảm ơn Thầy/Cô và các bạn đã lắng nghe!}}\\[0.5cm]
    \textit{Sẵn sàng trình diễn MVP và giải đáp câu hỏi.}
\end{center}
\end{frame}
